\section{\texorpdfstring{Example: a finite commutative ring, \(\mathbb{Z}/3\mathbb{Z}\)}{Example: a finite commutative ring, \textbackslash mathbb\{Z\}/3\textbackslash mathbb\{Z\}}}\label{sec:08-rings:05-finite-ring-example:1}

\texttt{intRing} is commutative, infinite, and --- being \(\mathbb{Z}\) itself --- the initial ring (Chapter 7's Theorem 1/2 hold for it as easily as possible, since it has essentially the ``obvious'' ring structure). Chapter 7's matrix example shows a genuinely \emph{noncommutative} ring. This section shows the other direction the axioms can specialize to: a \textbf{commutative, finite} ring, small enough to check every axiom by direct computation rather than by citing library lemmas.

\subsection{\texorpdfstring{The carrier: \texttt{Fin\ 3}}{The carrier: Fin 3}}\label{sec:08-rings:05-finite-ring-example:2}

\texttt{Fin 3} is Lean's type of ``naturals less than 3,'' i.e.~\(\{0, 1, 2\}\). Its built-in \texttt{+}/\texttt{*} already wrap around modulo 3, so \texttt{Fin 3}'s arithmetic \emph{is} \(\mathbb{Z}/3\mathbb{Z}\)'s arithmetic:

\begin{lstlisting}
#eval (2 : Fin 3) + 2   -- 1   (2 + 2 = 4 ≡ 1 mod 3)
#eval (2 : Fin 3) * 2   -- 1   (2 * 2 = 4 ≡ 1 mod 3)
\end{lstlisting}

\subsection{\texorpdfstring{Building \texttt{Group\ (Fin\ 3)}, \texttt{CommGroup\ (Fin\ 3)}, \texttt{Ring\ (Fin\ 3)}}{Building Group (Fin 3), CommGroup (Fin 3), Ring (Fin 3)}}\label{sec:08-rings:05-finite-ring-example:3}

Because \texttt{Fin 3} has only three elements, every one of \texttt{Group}/\texttt{Ring}'s axioms is a \textbf{finite, decidable} statement --- literally ``check this equation holds for all \(3\) (or \(3^2\), or \(3^3\)) choices of its variables.'' This is exactly the kind of goal Chapter 12 recommends handing to \href{https://lean-lang.org/doc/reference/latest/Tactic-Proofs/Tactic-Reference/}{\texttt{decide}} rather than proving by hand:

\begin{lstlisting}
def fin3Group : Group (Fin 3) where
  op := fun a b => a + b
  id := 0
  inv := fun a => -a
  assoc := by decide
  id_left := by decide
  id_right := by decide
  inv_left := by decide
  inv_right := by decide

def fin3CommGroup : CommGroup (Fin 3) where
  toGroup := fin3Group
  comm := by decide

def fin3Ring : Ring (Fin 3) where
  addGrp := fin3CommGroup
  mul := fun a b => a * b
  one := 1
  mul_assoc := by decide
  one_mul := by decide
  mul_one := by decide
  left_distrib := by decide
  right_distrib := by decide
\end{lstlisting}

\begin{progcorner}

\texttt{by decide} on a goal like \texttt{assoc} for \texttt{Fin 3} is not magic. It is exactly what one would get from

\end{progcorner}

\begin{lstlisting}[language=Python]
domain = range(3)
all((a + b) + c == a + (b + c) for a in domain for b in domain for c in domain)
\end{lstlisting}

brute-force enumerate every combination of the (finitely many) variables and check the equation holds for each. The difference is not the algorithm, but \emph{where} the check lives. \texttt{all(...)} is a runtime assertion: it only runs if some line of code calls it. If the call is omitted (or a later refactor deletes it), Python never notices that the property broke. \texttt{decide} instead runs as part of \emph{type-checking} \texttt{fin3Group} itself. The definition does not elaborate at all unless the brute-force check succeeds, so there is no such thing as a \texttt{fin3Group} value that quietly skipped its axiom checks. It is the same enumeration, but promoted from ``a test one hopes someone runs'' to a condition the kernel enforces before the value can even exist.

\begin{lstlisting}
#eval fin3Ring.addGrp.op 2 2          -- 1  (2 + 2 = 1 mod 3)
#eval fin3Ring.mul 2 2                 -- 1  (2 * 2 = 1 mod 3)
#eval fin3Ring.addGrp.toGroup.inv 1     -- 2  (-1 = 2 mod 3)
\end{lstlisting}

Compare this construction style directly with \texttt{intRing} (previous section) and \texttt{mat2Ring} (next section): \texttt{intRing}'s axioms needed \texttt{Int.add\textbackslash{}\_assoc}-style library citations because \texttt{Int} is infinite. There is no way to ``check all cases'' of a statement about \emph{every} integer. \texttt{Fin 3}'s axioms are statements about only finitely many (and concretely enumerable) elements, so they can instead be handled by exhaustive case-checking, which is exactly what \texttt{decide} automates. Recognizing which regime a goal is in --- infinite carrier needing a real argument, versus finite carrier where brute enumeration is enough --- is itself a useful instinct to build. See Chapter 12 for more on when \texttt{decide} is (and is not) the right tool.

\begin{mathreading}

\texttt{fin3Ring} is \(\mathbb{Z}/3\mathbb{Z}\), the finite field with three elements (though only the \texttt{Ring} structure has been built here --- a \texttt{Field} would additionally require every nonzero element to be invertible under multiplication, true for \(\mathbb{Z}/3\) precisely because \(3\) is prime, but this is not part of \texttt{Ring}'s axioms and not checked here). It is the quotient \(\mathbb{Z}/(3)\) by the ideal generated by \(3\), with \texttt{Fin 3}'s built-in wraparound arithmetic giving reduction mod \(3\) directly at the level of the representation, rather than through an explicit quotient construction.

\end{mathreading}

\textbf{Mathlib equivalent.} Mathlib's \href{https://loogle.lean-lang.org/?q=ZMod}{\texttt{ZMod 3}} \emph{is} \(\mathbb{Z}/3\mathbb{Z}\), already known to be a commutative ring (in fact a field, since \(3\) is prime). There is no \texttt{fin3Group}/\texttt{fin3CommGroup}/\texttt{fin3Ring} bundle, and no \texttt{decide} calls needed to re-verify axioms already proved once, generically, for every \texttt{n}:

\begin{lstlisting}
example : CommRing (ZMod 3) := inferInstance

#eval (2 : ZMod 3) + 2   -- 1
#eval (2 : ZMod 3) * 2   -- 1
\end{lstlisting}

The book's \texttt{by decide} on \texttt{fin3Ring}'s fields brute-forces exactly the finitely many cases needed for \emph{this one} carrier. Mathlib instead proves \texttt{ZMod n}'s ring axioms once for every \texttt{n} (through its construction as \texttt{Fin n} with wraparound arithmetic, combined with the general machinery for quotients of \(\mathbb{Z}\)), so \texttt{ZMod 3}'s instance costs nothing further to obtain.

\section{Next}\label{sec:08-rings:05-finite-ring-example:4}

Continue to \hyperref[sec:08-rings:06-accessing-fields:1]{Accessing nested fields}.
